% Provenance (never printed): see sources/outline_sources.md (Year 2: eccentricity, obliquity and
% precession; insolation computed from spherical astronomy; the Quaternary cycles; carbon dioxide
% and feedbacks). Orbital elements and tabulated insolation of Berger and Loutre (1991) via NOAA
% NCEI (data/derived/bachelor-2-orbit91, fetch_data Book 4 block); periods after Laskar et al.
% (2004) and the NASA Milankovitch page; history after the NASA Earth Observatory feature (2000);
% terminations after Barker and Knorr (2021); the mid-Pleistocene transition after Clark et al.
% (2006) and Koehler and van de Wal (2020). Figures from figdata/bachelor-2/orbital.py; values of
% the exercises and problem computed in the Book 4 scratch calculator calc/ch25.py.
\chapter{التأثير المداري والعصور الجليدية}\label{ch:b2:orbital-forcing}

على مدى آخر مليون سنة غيّر شكل مدار الأرض ضوء الشمس الذي يتلقّاه
الكوكب كلّه في سنة بأقلّ من جزأين في الألف: بين \num{340.00}
و\qty{340.54}{W.m^{-2}} في أعلى الغلاف الجوي. وعلى مدى المليون سنة نفسها تأرجح الإشماس اليومي المتوسط
في منتصف يوليو عند \ang{65} شمالًا بين نحو 388 و\qty{500}{W.m^{-2}}، أي بأكثر من
الربع. وقد تبعت العصور الجليدية الرقم الثاني، لا الأول. يحسب هذا الفصل كليهما
من عناصر مدار الأرض، ويسأل كيف يمكن لضوء الشمس عند عرض واحد في فصل واحد
أن يضبط إيقاع نموّ أغطية جليدية تغطّي قارّات وتلاشيها.
% ledger: data:b4-orbital, berger91:readme

\begin{recall}
\textbf{من كتب الفيزياء}: القطع الناقص، ونصف محوره الأكبر $a$ واختلافه المركزي $e$، وقوانين
كبلر. \textbf{من مجلد السنة الأولى}: دائرة البروج وميل المحور $\varepsilon$، وميل
الشمس $\delta$ مع $\sin\delta = \sin\varepsilon\sin\lambda$ (حيث $\lambda$ الطول الحقيقي للشمس،
المعدود على امتداد دائرة البروج من نقطة الاعتدال الربيعي)، والسنة المدارية، والإشماس
اليومي المتوسط
\[
  \bar Q = \frac{S_0}{\pi}\left(\frac{a}{r}\right)^{2}\bigl(h_0\sin\varphi\sin\delta +
  \cos\varphi\cos\delta\sin h_0\bigr),\qquad \cos h_0 = -\tan\varphi\tan\delta ,
\]
ومتوسطه العالمي السنوي، $S_0/(4\sqrt{1 - e^2})$. \textbf{من مجلّدات المدرسة}: الدورة
الجليدية والذروة الجليدية الأخيرة؛ والتأثير الإشعاعي. \textbf{من الفصول السابقة}:
\omterm{def:b2:precambrian-earth:ice-albedo}{التغذية الراجعة للجليد والألبيدو} (\cref{def:b2:precambrian-earth:ice-albedo})؛ والمكدّس القاعي، وسجلّ موقع EPICA
في القبّة C ونماذج العمر (\cref{ch:b2:climate-proxies}).
\end{recall}
% ledger: earth:TSI, earth:obliquity, earth:ecc

\section{المدار وعناصره الثلاثة}\label{sec:b2:orbital-forcing:elements}

مدار الأرض قطع ناقص تقع الشمس في إحدى بؤرتيه. والكواكب، والمشتري وزحل قبل
كلّ شيء، تشوّهه ببطء وتديره في مستويه؛ وجذب الشمس والقمر للانتفاخ
الاستوائي يجعل محور الدوران يترنّح كالنحلة الدوّارة. وثلاث كمّيات تصف ما يهمّ
ضوء الشمس: مدى استطالة القطع الناقص، ومدى ميل المحور، وموضع الفصول
على امتداد المدار.
% ledger: nasa:milank-causes

\begin{definition}[طول الحضيض؛ المبادرة المناخية]\label{def:b2:orbital-forcing:elements}
\emph{طول الحضيض}\index{طول الحضيض} $\varpi$ هو الزاوية، المقيسة في
مستوى المدار من اتجاه نقطة الاعتدال الربيعي، بين ذلك الاتجاه
واتجاه الحضيض كما يُرى من الشمس؛ وهو يُعدّ من نقطة الاعتدال المتحرّكة، فيتغيّر
بترنّح المحور وبدوران المدار معًا. و\emph{المبادرة
المناخية}\index{المبادرة المناخية} هي الكمّية $e\sin\varpi$، وتُسمّى أيضًا مؤشّر
المبادرة.
\end{definition}
% ledger: berger91:readme

\begin{omfigure}
\centering
\begin{tikzpicture}[scale=0.82, font=\scriptsize, >=Stealth]
\draw[thick, rotate=101.37] (-0.72,0) ellipse (2.4 and 2.2895);
\draw[densely dashed, black!60] (-0.331,1.647) -- (0.615,-3.059);
\draw[->, black!70] (0,0) -- (3.9,0);
\node[anchor=south, align=center] at (3.75,0.06) {\foreignlanguage{arabic}{$\gamma$، نحو}\\\foreignlanguage{arabic}{نقطة الاعتدال الربيعي}};
\draw[->, omDef] (0.75,0) arc[start angle=0, end angle=101.37, radius=0.75];
\node[omDef] at (0.78,0.62) {$\varpi$};
\fill[ocRed] (0,0) circle (0.13);
\node[anchor=north west] at (0.05,-0.08) {الشمس};
\foreach \x/\y in {-2.062/0, 0/-3.094, 2.321/0, 0/1.688} {\fill[ocBlue] (\x,\y) circle (0.08);}
\node[anchor=east, align=right] at (-2.15,0) {اعتدال\\مارس};
\node[anchor=north west, align=left] at (2.38,-0.08) {اعتدال\\سبتمبر};
\node[anchor=west] at (0.1,1.92) {انقلاب ديسمبر};
\node[anchor=east] at (-0.12,-3.25) {انقلاب يونيو};
\fill[black] (-0.331,1.647) circle (0.06);
\node[anchor=east] at (-0.42,1.72) {الحضيض};
\fill[black] (0.615,-3.059) circle (0.06);
\node[anchor=west] at (0.7,-3.15) {الأوج};
\draw[->, thick, black!60] (-1.9,-1.75) arc[start angle=220, end angle=245, radius=2.6];
\end{tikzpicture}
\par\smallskip{\small مدار الأرض كما يُرى من جهة الشمال، مع اختلاف مركزي مضخَّم
إلى 0.3 \omterm{def:b2:orbital-forcing:elements}{وطول الحضيض} الحالي، $\varpi = \ang{101.4}$. تتحرّك الأرض عكس اتجاه عقارب الساعة
(السهم)؛ وتحدّد النقاط موضعها عند الاعتدالين والانقلابين. ويقع الحضيض بعد بضعة
أيام من انقلاب ديسمبر، والأوج بعد بضعة أيام من انقلاب يونيو. تخطيطي.}
\end{omfigure}
% ledger: berger91:rows, nasa:milank-detail

كما تُرى من الأرض، تتحرّك الشمس على امتداد دائرة البروج، وطولها الحقيقي $\lambda$ هو
الزاوية من نقطة الاعتدال الربيعي: $\lambda = \ang{90}$ عند انقلاب يونيو، و\ang{270} عند انقلاب
ديسمبر. وتكون الأرض في الحضيض حين تكون الشمس في الاتجاه المقابل للحضيض،
عند $\lambda = \varpi + \ang{180}$؛ واليوم $\varpi = \ang{101.37}$ ويقع الحضيض عند
$\lambda = \ang{281.37}$، أي بعد نحو أحد عشر يومًا من انقلاب ديسمبر، في 3 يناير أو نحوه.
% ledger: berger91:rows, nasa:milank-detail

\begin{definition}[وسائط ميلانكوفيتش]\label{def:b2:orbital-forcing:milankovitch}
\emph{وسائط ميلانكوفيتش}\index{وسائط ميلانكوفيتش} هي الكمّيات المدارية الثلاث
التي تحدّد توزّع ضوء الشمس على الأرض عبر الفصول ودوائر
العرض: الاختلاف المركزي $e$، وميل المحور $\varepsilon$ \omterm{def:b2:orbital-forcing:elements}{والمبادرة المناخية}
$e\sin\varpi$. و\emph{نظرية ميلانكوفيتش}\index{نظرية ميلانكوفيتش} تنصّ على أن تغيّراتها
البطيئة تضبط إيقاع نموّ الأغطية الجليدية الرباعية وتلاشيها، عبر الإشماس الصيفي
في العروض الشمالية العليا.
\end{definition}
% ledger: nasa:milank-detail, eo00:milankovitch

لكلّ وسيط إيقاعاته. فللاختلاف المركزي أدوار رئيسية قرب \num{95000}
و\num{99000} و\num{124000} و\num{131000} سنة، تتجمّع في دورة من نحو
\num{100000} سنة، ودور قدره \num{405000} سنة، هو الأثبت بينها جميعًا؛ وعلى مدى آخر
مليون سنة تراوح بين \num{0.0046} و\num{0.0565}، وهو اليوم، عند \num{0.0172}، يتناقص
ببطء. ويتغيّر ميل المحور بدور رئيسي قدره \num{40978} سنة، بين \ang{22.1}
و\ang{24.5} على مدى آخر مليون سنة؛ وهو عند \ang{23.44} يتناقص، بعد أن تجاوز آخر قيمة عظمى له
قبل نحو \num{10000} سنة. ويترنّح المحور دورة كاملة في نحو \num{25770} سنة، بينما
يدور المدار في نحو \num{112000}؛ ومعًا يحرّكان الحضيض عبر الفصول بأدوار
رئيسية قرب \num{19000} و\num{22400} و\num{23700} سنة.
% ledger: la04:eccentricity, la04:obliquity, la04:precession, nasa:milank-detail, data:b4-orbital

\begin{omfigure}
\centering
% figdata/bachelor-2/orbital.py part elements
\begin{tikzpicture}
\begin{groupplot}[group style={group size=1 by 3, vertical sep=0.3cm, xlabels at=edge bottom,
  xticklabels at=edge bottom}, width=0.96\linewidth, height=3.1cm, xmin=0, xmax=1000,
  xtick={0,100,...,1000}, tick label style={font=\scriptsize}, label style={font=\small},
  scaled y ticks=false, yticklabel style={/pgf/number format/fixed,
  /pgf/number format/precision=2}]
\nextgroupplot[ylabel={$e$}, ymin=0, ymax=0.085, ytick={0,0.02,0.04,0.06}]
\addplot[omDef, thin] table[x=age_ka, y=ecc] {figdata/out/bachelor-2-orbital-elements.dat};
\node[font=\scriptsize, anchor=north west] at (rel axis cs:0.01,0.97) {أ\enspace الاختلاف المركزي};
\nextgroupplot[ylabel={\babelsublr{$\varepsilon$ (\unit{\degree})}}, ymin=21.9, ymax=25.4, ytick={22,23,24}]
\addplot[ocBlue, thin] table[x=age_ka, y=obl] {figdata/out/bachelor-2-orbital-elements.dat};
\node[font=\scriptsize, anchor=north west] at (rel axis cs:0.01,0.97) {ب\enspace ميل المحور};
\nextgroupplot[ylabel={$e\sin\varpi$}, ymin=-0.07, ymax=0.12, ytick={-0.05,0,0.05},
  xlabel={العمر بآلاف السنين قبل عام 1950}]
\addplot[ocRed, thin] table[x=age_ka, y=prec] {figdata/out/bachelor-2-orbital-elements.dat};
\node[font=\scriptsize, anchor=north west] at (rel axis cs:0.01,0.97) {ج\enspace المبادرة
المناخية};
\end{groupplot}
\end{tikzpicture}
\par\smallskip{\small \omterm{def:b2:orbital-forcing:milankovitch}{وسائط ميلانكوفيتش} الثلاثة على مدى آخر مليون سنة: (أ)
الاختلاف المركزي؛ (ب) ميل المحور؛ (ج) \omterm{def:b2:orbital-forcing:elements}{المبادرة المناخية}، التي غلافها الاختلاف المركزي.
المعطيات: بيرجيه ولوتر (1991)، قسم المناخ القديم في المراكز الوطنية للمعلومات البيئية التابعة للإدارة الوطنية للمحيطات والغلاف الجوي الأمريكية، كما جُلبت في 23 أكتوبر 2026.}
\end{omfigure}
% ledger: data:b4-orbit91

تتأرجح \omterm{def:b2:orbital-forcing:elements}{المبادرة المناخية} بسرعة، لكن بقدر ما يسمح لها الاختلاف المركزي فقط: فحين يكون المدار
شبه دائري، كما قبل \num{400000} سنة، لا يكاد يهمّ موضع الحضيض.
والشكل يُظهر ذلك بوضوح: غلاف اللوحة (ج) هو اللوحة (أ).
% ledger: data:b4-orbit91

\section{ضوء الشمس انطلاقًا من المدار}\label{sec:b2:orbital-forcing:insolation}

يحتاج الإشماس اليومي المتوسط المذكور أعلاه إلى عنصر آخر: بعد الشمس في
يوم معيّن.

\begin{proposition}[المسافة وطول الشمس]\label{prop:b2:orbital-forcing:distance}
حين يكون الطول الحقيقي للشمس $\lambda$، تكون المسافة بين الأرض والشمس
\[
  \frac{r}{a} = \frac{1 - e^2}{1 - e\cos(\lambda - \varpi)} .
\]
وعند انقلاب يونيو ($\lambda = \ang{90}$)، بتقريب من الرتبة الأولى في $e$،
$(a/r)^2 \approx 1 - 2e\sin\varpi$؛ وعند انقلاب ديسمبر، $(a/r)^2 \approx 1 + 2e\sin\varpi$.
\end{proposition}

\begin{proof}
على قطع ناقص تقع الشمس في بؤرته، $r/a = (1 - e^2)/(1 + e\cos v)$، حيث الشذوذ الحقيقي $v$
هو الزاوية عند الشمس من الحضيض إلى الأرض. والأرض تُرى من الشمس في
الاتجاه المقابل لاتجاه الشمس كما تُرى من الأرض، فطولها كما يُرى من الشمس هو
$\lambda + \ang{180}$، و$v = \lambda + \ang{180} - \varpi$، و$\cos v = -\cos(\lambda - \varpi)$.
وعند $\lambda = \ang{90}$، $\cos(\lambda - \varpi) = \sin\varpi$ و$(a/r)^2 = (1 -
e\sin\varpi)^2/(1 - e^2)^2 = 1 - 2e\sin\varpi + O(e^2)$؛ وعند $\lambda = \ang{270}$ تتغيّر الإشارة.
\end{proof}

\omterm{def:b2:orbital-forcing:elements}{فالمبادرة المناخية} إذن هي المؤشّر الصحيح: فهي تقيس مقدار قرب الأرض من الشمس، أو
بعدها عنها، عند انقلاب يونيو. فحين تكون $e\sin\varpi$ كبيرة موجبة، كما اليوم
وبقدر أكبر بكثير قبل \num{116000} سنة، يقع الصيف الشمالي قرب الأوج فيكون ضعيفًا؛ وحين تكون
كبيرة سالبة، كما قبل نحو \num{10000} سنة، يقع الصيف الشمالي قرب الحضيض.
% ledger: berger91:rows

\begin{example}[شمس يونيو عند \ang{65} شمالًا، اليوم]\label{ex:b2:orbital-forcing:today}
اليوم $e = 0.017236$، و$\varpi = \ang{101.37}$، و$\varepsilon = \ang{23.446}$، وتستعمل الجداول
$S_0 = \qty{1360}{W.m^{-2}}$. وعند انقلاب يونيو $\delta = \varepsilon$؛ وعند $\varphi =
\ang{65}$، $\cos h_0 = -\tan\ang{65}\tan\ang{23.446} = -0.9301$، فيكون $h_0 = \ang{158.44} =
\qty{2.7654}{rad}$: فالشمس فوق الأفق مدّة $2 \times 158.44/360 \times 24 = 21.1$~ساعة. وقوس
صيغة الإشماس هو $2.7654 \times 0.9063 \times 0.3979 + 0.4226 \times 0.9174 \times 0.3674 =
1.1397$، وعند المسافة المتوسطة $\bar Q = 432.90 \times 1.1397 = \qty{493.4}{W.m^{-2}}$.
وعامل المسافة $(a/r)^2 = 0.9671$، فيكون $\bar Q = \qty{477.1}{W.m^{-2}}$.
\end{example}
% ledger: berger91:rows, berger91:readme, data:b4-orbital

يضيف قانون كبلر الثاني لمسة أخرى. فقرب الأوج تتحرّك الأرض أبطأ: فالفصل الذي يقع
هناك أضعف لكنه أطول. والصيف الشمالي اليوم أطول بنحو 4.5 يوم من الشتاء
الشمالي، والربيع أطول بنحو ثلاثة أيام من الخريف.
% ledger: nasa:milank-detail

\begin{proposition}[المتوسطات السنوية لا تتوقّف على المبادرة]\label{prop:b2:orbital-forcing:annual}
الطاقة المتلقّاة لكلّ وحدة مساحة في أعلى الغلاف الجوي بين طولين معيّنين
للشمس لا تتوقّف على $\varpi$، والإشماس السنوي المتوسط عند أيّ دائرة عرض هو إشماس
مدار دائري نصف قطره $a$ مضروبًا في $1/\sqrt{1 - e^2}$. وعند القطب يكون المتوسط السنوي
\[
  \bar Q_{\text{القطب}} = \frac{S_0\sin\varepsilon}{\pi\sqrt{1 - e^2}} ,
\]
وهو يزداد مع ميل المحور.
\end{proposition}

\begin{proof}
بقانون كبلر الثاني يمسح نصف القطر المساحة بالمعدّل الثابت $\pi ab/T$، والمساحة الممسوحة
بينما تدور الأرض بمقدار $\mathrm d\lambda$ هي $\frac12 r^2\,\mathrm d\lambda$، فيكون $\mathrm dt =
T r^2\,\mathrm d\lambda/(2\pi ab)$، مع $b = a\sqrt{1 - e^2}$. والإشماس عند عرض معيّن
وطول معيّن للشمس هو $(a/r)^2$ مضروبًا في قيمته $Q_a(\lambda)$ عند المسافة $a$، فتكون
الطاقة المتلقّاة $\int Q_a(\lambda)(a/r)^2\,\mathrm dt = \frac{T}{2\pi\sqrt{1 - e^2}}\int
Q_a(\lambda)\,\mathrm d\lambda$: فقد اختزلت المسافة، ومعها $\varpi$. وعلى مدى السنة
كلّها، يجري $\lambda$ عبر \ang{360} ويكون المتوسط $\frac{1}{2\pi\sqrt{1-e^2}}\int_0^{2\pi}
Q_a\,\mathrm d\lambda$. وعند القطب الشمالي $Q_a = S_0\sin\delta = S_0\sin\varepsilon\sin\lambda$
مع $0 < \lambda < \pi$ و0 فيما عدا ذلك؛ والتكامل يساوي $2S_0\sin\varepsilon$.
\end{proof}
% ledger: earth:TSI

المبادرة تنقل ضوء الشمس من فصل إلى آخر، ولا تُدخله في السنة أو تُخرجه منها أبدًا؛ والاختلاف المركزي
يغيّر المجموع السنوي بالعامل $1/\sqrt{1 - e^2}$، أي بنسبة \qty{0.16}{\percent} على الأكثر على مدى
آخر مليون سنة؛ وميل المحور ينقل ضوء الشمس بين دوائر العرض، فيرفع المتوسط السنوي عند
القطبين ويخفضه قليلًا عند خطّ الاستواء. ولا يتغيّر كثيرًا إلا النمط الفصلي والعرضي،
وهذا ما يبيّنه الشكل أدناه.
% ledger: data:b4-orbital, nasa:milank-detail

\begin{omfigure}
\centering
% figdata/bachelor-2/orbital.py part lat
\begin{tikzpicture}
\begin{groupplot}[group style={group size=2 by 1, horizontal sep=1.6cm},
  tick label style={font=\scriptsize}, label style={font=\small}, height=5.4cm,
  width=0.48\linewidth, xmin=-90, xmax=90, xtick={-90,-60,...,90}, xlabel={العرض الجغرافي
  \babelsublr{(\unit{\degree})}}, ylabel={\babelsublr{$\bar Q$ (\unit{W.m^{-2}})}}]
\nextgroupplot[ymin=0, ymax=650, title={أ\enspace انقلاب يونيو}, title style={font=\scriptsize},
  legend style={font=\scriptsize, at={(0.03,0.97)}, anchor=north west, draw=none},
  legend cell align=left]
\addplot[black, thick] table[x=lat, y=jun_0] {figdata/out/bachelor-2-orbital-lat.dat};
\addlegendentry{اليوم}
\addplot[ocRed, thick, densely dashed] table[x=lat, y=jun_10] {figdata/out/bachelor-2-orbital-lat.dat};
\addlegendentry{\foreignlanguage{arabic}{قبل \num{10000} سنة}}
\nextgroupplot[ymin=100, ymax=500, title={ب\enspace المتوسط السنوي}, title style={font=\scriptsize},
  legend style={font=\scriptsize, at={(0.5,0.03)}, anchor=south, draw=none},
  legend cell align=left]
\addplot[ocBlue, thick] table[x=lat, y=ann_lo] {figdata/out/bachelor-2-orbital-lat.dat};
\addlegendentry{$\varepsilon = \ang{22.1}$}
\addplot[ocRed, thick, densely dashed] table[x=lat, y=ann_hi] {figdata/out/bachelor-2-orbital-lat.dat};
\addlegendentry{$\varepsilon = \ang{24.5}$}
\end{groupplot}
\end{tikzpicture}
\par\smallskip{\small (أ) الإشماس اليومي المتوسط عند انقلاب يونيو، اليوم وقبل \num{10000} سنة،
حين كان الصيف الشمالي يقع قرب الحضيض وكان ميل المحور أكبر: أي نحو
\qty{50}{W.m^{-2}} زيادةً في العروض الشمالية العليا. (ب) الإشماس السنوي المتوسط لقيمتي
ميل المحور القصويين في آخر مليون سنة، مع الاختلاف المركزي الحالي: فالميل الأكبر يرسل
مزيدًا من ضوء الشمس إلى القطبين وأقلّ قليلًا إلى المنطقة المدارية. محسوب من العناصر المدارية
مع $S_0 = \qty{1360}{W.m^{-2}}$.}
\end{omfigure}
% ledger: data:b4-orbital, data:b4-orbit91

\section{حساب منحنى الإشماس}\label{sec:b2:orbital-forcing:curve}

صار منحنى الإشماس الذي بناه ميلانكوفيتش اليوم حسابًا روتينيًا. فالعناصر المدارية
للماضي تُحسب بمكاملة حركات الكواكب؛ وهي، مجدولةً كلّ ألف سنة،
المدخلات. ولكلّ عمر، تعطي الطريقة أدناه الإشماس عند أيّ عرض وفي أيّ فصل.

\begin{method}[إشماس الماضي]\label{met:b2:orbital-forcing:curve}
\begin{enumerate}
  \item اقرأ $e$ و$\varpi$ و$\varepsilon$ للعمر المطلوب.
  \item اختر الفصل بالطول الحقيقي للشمس $\lambda$: \ang{90} لانقلاب يونيو؛
        ومنتصف يوليو في الجداول يقابل $\lambda = \ang{120}$.
  \item احسب $\delta$ من $\sin\delta = \sin\varepsilon\sin\lambda$، والزاوية الساعية
        للغروب $h_0$ عند العرض $\varphi$ (نهار قطبي أو ليل قطبي إذا كان $|\tan\varphi\tan\delta| > 1$).
  \item احسب عامل المسافة $(a/r)^2$ بواسطة \cref{prop:b2:orbital-forcing:distance}.
  \item اضرب: $\bar Q = (S_0/\pi)(a/r)^2(h_0\sin\varphi\sin\delta +
        \cos\varphi\cos\delta\sin h_0)$.
\end{enumerate}
\end{method}
% ledger: berger91:readme, data:b4-orbital

إذا طُبّقت الطريقة على أعمار آخر مليون سنة البالغ عددها 1001، فإنها تعيد إنتاج الإشماس المجدول لمنتصف يوليو
عند \ang{65} شمالًا بفارق لا يتجاوز \qty{0.012}{W.m^{-2}}، وهو تقريب الجداول،
وقيم يونيو المجدولة لليوم الحاضر تمامًا. وتضع الجداول مبدأ الزمن عند عام 1950.
% ledger: data:b4-orbital, berger91:readme

\begin{omfigure}
\centering
% figdata/bachelor-2/orbital.py part jul65
\begin{tikzpicture}
\begin{groupplot}[group style={group size=1 by 2, vertical sep=0.3cm, xlabels at=edge bottom,
  xticklabels at=edge bottom}, width=0.96\linewidth, xmin=0, xmax=400, xtick={0,50,...,400},
  tick label style={font=\scriptsize}, label style={font=\small}]
\nextgroupplot[height=5.0cm, ymin=370, ymax=590, ylabel={$\bar Q$ عند \babelsublr{\ang{65}} شمالًا
  \babelsublr{(\unit{W.m^{-2}})}}, legend style={font=\scriptsize, at={(0.5,1.02)}, anchor=south, draw=none,
  legend columns=3, /tikz/every even column/.append style={column sep=6pt}},
  legend cell align=left]
\addplot[black!45, line width=2.2pt] table[x=age_ka, y=tab] {figdata/out/bachelor-2-orbital-jul65.dat};
\addlegendentry{منتصف يوليو، مجدول}
\addplot[ocRed, thin] table[x=age_ka, y=calc] {figdata/out/bachelor-2-orbital-jul65.dat};
\addlegendentry{منتصف يوليو، محسوب}
\addplot[ocBlue, thin, densely dashed] table[x=age_ka, y=jun65] {figdata/out/bachelor-2-orbital-jul65.dat};
\addlegendentry{انقلاب يونيو، محسوب}
\nextgroupplot[height=2.8cm, ymin=-0.02, ymax=0.02, ytick={-0.01,0,0.01},
  scaled y ticks=false, yticklabel style={/pgf/number format/fixed}, ylabel={المتبقّي},
  xlabel={العمر بآلاف السنين قبل عام 1950}]
\addplot[only marks, mark=*, mark size=0.35pt, black] table[x=age_ka, y=resid]
  {figdata/out/bachelor-2-orbital-jul65.dat};
\end{groupplot}
\end{tikzpicture}
\par\smallskip{\small في الأعلى: الإشماس اليومي المتوسط عند \ang{65} شمالًا خلال آخر \num{400000} سنة،
لمنتصف يوليو كما جدوله بيرجيه ولوتر (رمادي، عريض) وكما حُسب بالطريقة أعلاه (أحمر)،
وعند انقلاب يونيو (أزرق، متقطّع). في الأسفل: القيم المحسوبة ناقص القيم المجدولة لمنتصف يوليو،
بوحدة \unit{W.m^{-2}}.}
\end{omfigure}
% ledger: data:b4-orbit91, data:b4-orbital

تبرز سمتان. فالمنحنى تهيمن عليه \omterm{def:b2:orbital-forcing:elements}{المبادرة المناخية}، بقمم يفصل بينها نحو
\num{20000} سنة يتضخّم ارتفاعها ويتقلّص مع الاختلاف المركزي؛ ويضيف ميل المحور
تعديلًا أبطأ. والتأرجحات كبيرة: فخلال آخر \num{800000} سنة تراوحت قيمة منتصف يوليو
بين 388 و\qty{497}{W.m^{-2}}.
% ledger: data:b4-orbital

\section{نظرية ميلانكوفيتش والدورات الجليدية}\label{sec:b2:orbital-forcing:theory}

\begin{history}[ضوء الشمس عند خمس وستين درجة شمالًا]
بنى ميلوتين ميلانكوفيتش، المولود عام 1879، نظرية رياضية للمناخ انطلاقًا من التغيّرات الفصلية والعرضية لضوء الشمس. وبجمع التغيّرات المدارية الثلاثة،
حسب إشماس دوائر العرض ودرجات حرارة السطح المقابلة لها على مدى
\num{600000} سنة قبل عام 1800. وباقتراح من عالم المناخ كوبن اختار الإشماس
الصيفي عند \ang{65} شمالًا كمّيةً حاسمة: فالأغطية الجليدية الكبرى نمت قرب تلك
الدائرة، والصيف البارد يُذيب قدرًا أقلّ من ثلج الشتاء، فيتراكم الثلج من سنة إلى
أخرى. وقد عدّ ميل المحور أهمّ الوسائط الثلاثة. وطوال نحو خمسين سنة
بقيت النظرية مهمَلة إلى حدّ بعيد؛ وتوفّي ميلانكوفيتش عام 1958. وفي عام 1976 وجد هايز وإمبري وشاكلتون
في رواسب أعماق البحار لآخر \num{450000} سنة قممًا طيفية لتقلّبية المناخ
عند نحو \num{23000} و\num{42000} و\num{100000} سنة، وهي أدوار المبادرة وميل المحور
والاختلاف المركزي، واستنتجوا أن تغيّرات مدار الأرض هي السبب الأساسي
للعصور الجليدية الرباعية.
\end{history}
% ledger: eo00:milankovitch, nasa:milank-detail, hays76:periods

يكرّر الشكل أدناه اختبارهم على سجلّات \cref{ch:b2:climate-proxies}. وهو يقارن
أطياف السعة، أي شدّة كلّ دور، لإشماس يوليو عند \ang{65} شمالًا وللمكدّس
القاعي، في نافذتين طول كلّ منهما \num{700000} سنة.

\begin{omfigure}
\centering
% figdata/bachelor-2/orbital.py part spectrum
\begin{tikzpicture}
\begin{axis}[width=0.9\linewidth, height=5.6cm, xmin=10, xmax=150, ymin=0, ymax=1.08,
  xlabel={الدور بآلاف السنين}, ylabel={السعة النسبية},
  xtick={19,23,41,100,150}, tick label style={font=\scriptsize}, label style={font=\small},
  legend style={font=\scriptsize, at={(0.5,1.02)}, anchor=south, draw=none, legend columns=3,
  /tikz/every even column/.append style={column sep=6pt}}, legend cell align=left]
\addplot[black!50, thick] table[x=period, y=amp_ins] {figdata/out/bachelor-2-orbital-spectrum.dat};
\addlegendentry{\foreignlanguage{arabic}{الإشماس عند \ang{65} شمالًا، \babelsublr{0--\qty{700}{ka}}}}
\addplot[ocRed, thick] table[x=period, y=amp_young] {figdata/out/bachelor-2-orbital-spectrum.dat};
\addlegendentry{\babelsublr{LR04, 0--\qty{700}{ka}}}
\addplot[ocBlue, thick, densely dashed] table[x=period, y=amp_old] {figdata/out/bachelor-2-orbital-spectrum.dat};
\addlegendentry{\babelsublr{LR04, \numrange{1250}{1950}~ka}}
\end{axis}
\end{tikzpicture}
\par\smallskip{\small أطياف السعة، وكلّ منها مقيس إلى قيمته العظمى، لإشماس منتصف يوليو
عند \ang{65} شمالًا خلال آخر \num{700000} سنة، ولمكدّس LR04 القاعي خلال النافذة نفسها
وخلال الفترة من \num{1250000} إلى \num{1950000} سنة مضت (بعد إزالة الاتجاه العام وتطبيق نافذة تخفيف). ويتركّز
نشاط الإشماس عند أدوار المبادرة، قرب 19 إلى 24 ألف سنة، وعند دور ميل المحور؛ أمّا
حجم الجليد الحديث فيبلغ ذروته قرب \num{100000} سنة، حيث لا يكاد يوجد في الإشماس شيء؛ وحجم الجليد
الأقدم يبلغ ذروته عند \num{41000} سنة.}
\end{omfigure}
% ledger: data:b4-orbital, data:lr04, hays76:periods

\subsection*{الفترات الجليدية والفترات بين الجليدية ونهايات الفترات الجليدية}

\begin{definition}[الفترة الجليدية؛ الفترة بين الجليدية؛ نهاية الفترة الجليدية]\label{def:b2:orbital-forcing:glacials}
في سجلّ الرباعي، \emph{الفترة الجليدية}\index{الفترة الجليدية} مجال زمني تغطّي فيه أغطية
جليدية كبيرة القارّات الشمالية، ويُسجَّل بقيم مرتفعة للمؤشّر القاعي \mbox{\textdir TLT $\delta^{18}$O} وبمستوى منخفض للبحر؛
و\emph{الفترة بين الجليدية}\index{الفترة بين الجليدية} مجال زمني قليل الجليد البرّي في نصف الكرة
الشمالي خارج غرينلاند، كالحاضر. و\emph{نهاية الفترة
الجليدية}\index{نهاية الفترة الجليدية} واحد من الانتقالات الكبرى التي تُنهي دورة جليدية،
من الظروف الجليدية التامّة إلى ظروف بين جليدية، وهو أسرع بكثير من النموّ الذي سبقه.
\end{definition}
% ledger: kohler20:mpt, barker21:terminations, ipcc:ref-periods

خلال آخر نصف مليون سنة يقترب الدور المهيمن للمناخ من \num{100000} سنة،
مع أن التأثير عند ذلك الدور، وهو آتٍ من الاختلاف المركزي وحده، مهمل
إلى جانب تأثير ميل المحور والمبادرة. والدورات غير متناظرة أيضًا: فالمؤشّر القاعي \mbox{\textdir TLT $\delta^{18}$O}
يسجّل تراكمًا متدرّجًا للجليد على مدى نحو \num{90000} سنة، ثم اضمحلالًا في نحو
\num{10000}، أي سنًّا منشاريًا حوافّه الهابطة هي نهايات الفترات الجليدية. وآخر فترة بين جليدية، من
نحو \num{129000} إلى \num{116000} سنة مضت، انتهت حين هبط الإشماس الصيفي عند \ang{65} شمالًا
إلى واحدة من أدنى قيمه في آخر \num{800000} سنة: فقد بلغت قيمة انقلاب يونيو
أدناها، \qty{438.6}{W.m^{-2}}، قبل \num{116000} سنة، وبلغت قيمة منتصف يوليو ثاني أدنى
قيمة دنيا لها، \qty{392.2}{W.m^{-2}}، قبل \num{114000} سنة. وجاءت الذروة الجليدية الأخيرة بعد نحو
\num{95000} سنة.
% ledger: barker21:terminations, ipcc:ref-periods, data:b4-orbital, data:b4-orbital-extremes

ما مدى سرعة غطاء جليدي في متابعة تأثيره؟ يعامل أبسط نموذج حجم الجليد $V$ كأنه
يسترخي نحو الحجم $V_{\mathrm{eq}}$ الذي يكفي الإشماس الحالي لإبقائه، مع
زمن استجابة $\tau$:
\[
  \frac{\mathrm dV}{\mathrm dt} = \frac{V_{\mathrm{eq}}(t) - V}{\tau} .
\]

\begin{proposition}[استجابة متأخّرة]\label{prop:b2:orbital-forcing:lag}
لتأثير $V_{\mathrm{eq}} = A\sin\omega t$ يؤول حلّ نموذج التأخّر إلى
\[
  V = \frac{A}{\sqrt{1 + \omega^2\tau^2}}\,\sin\bigl(\omega(t - t_{\ell})\bigr),\qquad
  t_{\ell} = \frac{\arctan(\omega\tau)}{\omega} :
\]
فالجليد يتابع التأثير البطيء متابعةً شبه تامّة والتأثير السريع متابعةً ضعيفة، ودائمًا متأخّرًا. والاستجابة
لمجموع من الدوال الجيبية هي مجموع الاستجابات: فالنموذج لا يُنشئ دورًا يفتقر إليه
تأثيره.
\end{proposition}

\begin{proof}
لنجرّب $V = B\sin\omega t + C\cos\omega t$: بالتعويض، $\omega B\tau = -C$ و$-\omega C\tau = A -
B$، فيكون $B = A/(1 + \omega^2\tau^2)$، و$C = -\omega\tau A/(1 + \omega^2\tau^2)$، أي سعة
$\sqrt{B^2 + C^2} = A/\sqrt{1 + \omega^2\tau^2}$ وتأخّر في الطور قدره $\arctan(\omega\tau)$. وأيّ حلّ
آخر يختلف عن هذا بمضاعف للدالة $\mathrm e^{-t/\tau}$ التي تتلاشى.
والمعادلة خطّية، فالحلول تُجمع.
\end{proof}

\begin{omfigure}
\centering
% figdata/bachelor-2/orbital.py part lag
\begin{tikzpicture}
\begin{axis}[width=0.96\linewidth, height=5.4cm, xmin=0, xmax=400, ymin=-3.4, ymax=4.2,
  xtick={0,50,...,400}, xlabel={العمر بآلاف السنين قبل عام 1950},
  ylabel={قيم معيارية}, tick label style={font=\scriptsize}, label style={font=\small},
  legend style={font=\scriptsize, at={(0.5,1.02)}, anchor=south, draw=none, legend columns=3,
  /tikz/every even column/.append style={column sep=6pt}}, legend cell align=left]
\addplot[black!35, thin] table[x=age_ka, y=forcing] {figdata/out/bachelor-2-orbital-lag.dat};
\addlegendentry{\foreignlanguage{arabic}{سالب الإشماس عند \ang{65} شمالًا}}
\addplot[ocRed, thick] table[x=age_ka, y=model] {figdata/out/bachelor-2-orbital-lag.dat};
\addlegendentry{\foreignlanguage{arabic}{نموذج التأخّر، \babelsublr{$\tau = \qty{10}{kyr}$}}}
\addplot[black, thick] table[x=age_ka, y=lr04] {figdata/out/bachelor-2-orbital-lag.dat};
\addlegendentry{\foreignlanguage{arabic}{\babelsublr{LR04} (الجليد إلى الأعلى)}}
\end{axis}
\end{tikzpicture}
\par\smallskip{\small نموذج التأخّر مدفوعًا بإشماس منتصف يوليو عند \ang{65} شمالًا (ضوء
شمس أقلّ، جليد أكثر)، مقابل المكدّس القاعي، والثلاثة بقيم معيارية (متوسط 0، وانحراف معياري
1)؛ والقيم الموجبة تعني جليدًا أكثر. يتابع النموذج تموّجات المبادرة في المكدّس لكنه يخلو
من دوراته المنشارية الكبيرة. وزمن الاستجابة البالغ \qty{10}{kyr} اختيار في النموذج.}
\end{omfigure}
% ledger: data:b4-orbital, data:lr04

خلال آخر \num{800000} سنة يرتبط النموذج المتأخّر بالمكدّس بمعامل $r = 0.42$، أمّا
الإشماس نفسه فبمعامل $0.23$ فقط: فالتأخّر يساعد، لكن غطاءً جليديًا خطّيًا لا يستطيع أن يصنع
دورات من \num{100000} سنة من تأثير يخلو منها. فلا بدّ أن تنطوي نهايات الفترات الجليدية على
سلوك غير خطّي يضخّم دفعة متواضعة من الإشماس. والمرشّحون تغذيات راجعة تؤثّر في المناخ كلّه:
\omterm{def:b2:precambrian-earth:ice-albedo}{التغذية الراجعة للجليد والألبيدو} في \cref{def:b2:precambrian-earth:ice-albedo}، لأن كوكبًا فيه مزيد
من الثلج والجليد يعكس مزيدًا من ضوء الشمس؛ وثنائي أكسيد الكربون، الذي يرتفع في لبّات الجليد القطبية الجنوبية على
مقاييس زمنية من آلاف السنين كلّما ضعفت الدورة الانقلابية في الأطلسي، كما
يحدث حين تتدفّق مياه ذوبان الأغطية الجليدية في شمال الأطلسي. وبعد عتبة من ثنائي أكسيد
الكربون، يمكن للظروف بين الجليدية أن تستقرّ.
% ledger: data:b4-orbital, barker21:terminations, eo00:milankovitch

\begin{definition}[الانتقال في منتصف البليستوسين]\label{def:b2:orbital-forcing:mpt}
\emph{الانتقال في منتصف البليستوسين}\index{الانتقال في منتصف البليستوسين} هو التحوّل، من نحو
\num{1250000} إلى \num{700000} سنة مضت، من دورات جليدية من نحو \num{41000} سنة وسعة
صغيرة إلى دورات شبه دورية من نحو \num{100000} سنة وسعة كبيرة، مع أغطية جليدية
أسمك، من دون أيّ تغيّر مقابل في التأثير المداري.
\end{definition}
% ledger: clark06:mpt, nasa:milank-detail, kohler20:mpt

قبل الانتقال كانت الأغطية الجليدية تتبع ميل المحور، كما يُظهر الطيف المتقطّع في الشكل
أعلاه؛ ويجد أحد الإحصاءات نهايةً لفترة جليدية في 67 و88 و\qty{52}{\percent} من
دورات ميل المحور في الرباعي المبكّر والأوسط والمتأخّر. والنظرية المدارية وحدها لا تستطيع تفسير هذا التحوّل؛ والمرشّحون هم ديناميكية
الأغطية الجليدية نفسها، ربّما بعد إزالة الحطام الصخري الليّن الذي كانت تنزلق عليه،
وتراجع تدريجي لثنائي أكسيد الكربون خلال الحدود الدنيا الجليدية. ولا يوجد بعدُ تفسير
متّفق عليه عمومًا.
% ledger: kohler20:mpt, nasa:milank-detail, clark06:mpt

\section{ثنائي أكسيد الكربون في الدورات الجليدية}\label{sec:b2:orbital-forcing:co2}

تُظهر لبّات الجليد القطبية الجنوبية أن ثنائي أكسيد الكربون يتبع الدورات الجليدية: بين 188
و\qty{194}{ppm} في الذروة الجليدية الأخيرة، وبين 266 و\qty{282}{ppm} في آخر فترة بين جليدية؛
ويعطي السجلّ المركّب في \cref{ch:b2:climate-proxies} القيمة \qty{189.7}{ppm} للفترة من 19 إلى
\qty{23}{ka} والقيمة \qty{279.4}{ppm} على مدى الألفي سنة السابقة للصناعة. ويزداد تأثيره الإشعاعي
بما يقارب لوغاريتم التركيز، أي \qty{3.93}{W.m^{-2}} لكلّ مضاعفة.
% ledger: ipcc:fig2-34, data:b4-proxies-co2, ipcc:log-co2, ipcc:erf-2xco2

\begin{example}[المدار مقابل ثنائي أكسيد الكربون]\label{ex:b2:orbital-forcing:co2}
عند الانتقال من \qty{189.7}{ppm} إلى \qty{279.4}{ppm}، يكون التأثير الإشعاعي $3.93\log_2(279.4/189.7) = 3.93 \times
0.559 = \qty{2.2}{W.m^{-2}}$، على الكرة الأرضية كلّها وعلى مدى السنة كلّها. أمّا المدار، على مدى آخر
مليون سنة، فقد غيّر المتوسط العالمي السنوي للإشماس بما لا يزيد على $340.54 - 340.00 =
\qty{0.54}{W.m^{-2}}$، لكنه غيّر إشماس منتصف يوليو عند \ang{65} شمالًا بنحو \qty{109}{W.m^{-2}}.
\end{example}
% ledger: data:b4-proxies-co2, ipcc:erf-2xco2, ipcc:log-co2, data:b4-orbital

المدار ضابط إيقاع، لا محرّك. فهو يعيد توزيع ضوء الشمس بين الفصول ودوائر العرض
ويحدّد التوقيت؛ أمّا الأغطية الجليدية، بألبيدوها وبطئها، وثنائي أكسيد الكربون، بتأثيره
العالمي، فيحوّلان إشارة إقليمية إلى تغيّر عالمي في المناخ.
% ledger: hays76:periods, barker21:terminations, eo00:milankovitch

\section*{التمارين}\label{sec:b2:orbital-forcing:exercises}

معطيات التمارين، ما لم يُذكر غير ذلك: $S_0 = \qty{1360}{W.m^{-2}}$ (قيمة
الجداول)؛ وسنة مدارية من \num{365.2422} يومًا؛ والعناصر المدارية في الجداول، وقيمها اليوم
$e = 0.017236$، و$\varpi = \ang{101.37}$، و$\varepsilon = \ang{23.446}$.
% ledger: berger91:readme, berger91:rows

\begin{exercise}[$\star$]\label{exo:b2:orbital-forcing:1}
خذ $e = 0.0167$. (أ) احسب نسبة مسافة الأوج $a(1 + e)$ إلى مسافة
الحضيض $a(1 - e)$. (ب) عبّر عن الفرق بينهما نسبةً مئوية من $a$. (ج) بأيّ عامل يكون
ضوء الشمس عند الحضيض أقوى منه عند الأوج؟
\end{exercise}
% ledger: earth:ecc, nasa:milank-detail

\begin{exercise}[$\star$]\label{exo:b2:orbital-forcing:2}
تعطي الجداول $e = 0.017236$، و$\varpi = \ang{101.37}$ اليوم؛ و$e = 0.019971$، و$\varpi =
\ang{293.86}$ قبل \num{10000} سنة؛ و$e = 0.019249$، و$\varpi = \ang{359.99}$ قبل \num{6000} سنة.
(أ) احسب \omterm{def:b2:orbital-forcing:elements}{المبادرة المناخية} لكلّ عمر. (ب) أعطِ لكلّ منها طول الشمس عند
الحضيض. (ج) في أيّ عمر كان الحضيض أقرب ما يكون إلى انقلاب يونيو، وفي أيّ عمر إلى
اعتدال سبتمبر؟
\end{exercise}
% ledger: berger91:rows

\begin{exercise}[$\star$]\label{exo:b2:orbital-forcing:3}
يتحرّك محور الدوران حركة مبادرة بدور قدره نحو \num{25770} سنة بالنسبة إلى النجوم، بينما
يدور المدار بدور قدره نحو \num{112000} سنة. افترض أن الحركتين كلتيهما تحرّكان
الحضيض بالنسبة إلى نقطة الاعتدال في الاتجاه نفسه، فيُجمع تردّداهما. (أ) اجمع
التردّدين لإيجاد الدور الذي يتحرّك به الحضيض
عبر الفصول. (ب) قارن بالأدوار الرئيسية للمبادرة المناخية، \num{19.0}
و\num{22.4} و\num{23.7} ألف سنة.
\end{exercise}
% ledger: nasa:milank-detail, la04:precession

\begin{exercise}[$\star$]\label{exo:b2:orbital-forcing:4}
باستعمال \cref{prop:b2:orbital-forcing:annual} و$e = 0.0167$، احسب الإشماس السنوي المتوسط عند
القطب الشمالي لقيمتي ميل المحور \ang{22.1} و\ang{24.5}. بأيّ نسبة مئوية يزداد بين
القيمتين؟
\end{exercise}
% ledger: nasa:milank-detail, earth:ecc

\begin{exercise}[$\star\star$]\label{exo:b2:orbital-forcing:5}
(أ) احسب بدقّة عامل المسافة $(a/r)^2$ عند انقلاب يونيو اليوم وقبل \num{116000} سنة
($e = 0.043988$، $\varpi = \ang{92.71}$)، وقارن بالتقريب $1 -
2e\sin\varpi$. (ب) احسب كليهما عند انقلاب ديسمبر. (ج) لماذا يكون التقريب أسوأ قبل \num{116000}
سنة؟
\end{exercise}
% ledger: berger91:rows

\begin{exercise}[$\star\star$]\label{exo:b2:orbital-forcing:6}
أعِد إنتاج \cref{ex:b2:orbital-forcing:today} خطوة خطوة: الزاوية الساعية للغروب عند \ang{65} شمالًا
عند انقلاب يونيو، وطول النهار، والإشماس عند المسافة المتوسطة،
وإشماس اليوم.
\end{exercise}
% ledger: berger91:rows

\begin{exercise}[$\star\star$]\label{exo:b2:orbital-forcing:7}
(أ) باستعمال $\mathrm dt = T(r/a)^2\,\mathrm d\lambda/(2\pi\sqrt{1 - e^2})$ و$(r/a)^2 \approx 1 +
2e\cos(\lambda - \varpi)$، بيّن أن نصف السنة من اعتدال مارس إلى اعتدال سبتمبر يدوم
نحو $T/2 + (2T/\pi)\,e\sin\varpi$. (ب) احسب بكم يومًا يزيد على نصف السنة الآخر
اليوم. (ج) قارن بمدّة 4.5 يوم التي يزيد بها الصيف الشمالي على الشتاء وبالأيام الثلاثة التي
يزيد بها الربيع على الخريف.
\end{exercise}
% ledger: nasa:milank-detail, berger91:rows

\begin{exercise}[$\star\star$]\label{exo:b2:orbital-forcing:8}
خلال آخر مليون سنة تراوح الاختلاف المركزي بين \num{0.0046} و\num{0.0565}. (أ) احسب
الإشماس العالمي السنوي المتوسط $S_0/(4\sqrt{1 - e^2})$ عند الطرفين. (ب) بكم غيّره المدار،
بوحدة \unit{W.m^{-2}} وبالنسبة المئوية؟ (ج) قارن بمدى إشماس
منتصف يوليو عند \ang{65} شمالًا، من 388 إلى \qty{497}{W.m^{-2}}.
\end{exercise}
% ledger: data:b4-orbital

\begin{exercise}[$\star\star$]\label{exo:b2:orbital-forcing:9}
باستعمال نموذج التأخّر و$\tau = \qty{10}{kyr}$، احسب السعة النسبية والتأخّر لحجم
الجليد لتأثير جيبي أدواره 23 و41 و\num{100} ألف سنة. أيّ نطاق يفضّله
الغطاء الجليدي، ولماذا؟
\end{exercise}

\begin{exercise}[$\star\star\star$]\label{exo:b2:orbital-forcing:10}
إشماس منتصف يوليو عند \ang{65} شمالًا لا يكاد يملك قدرة عند \num{100000} سنة، ومع ذلك فإن حجم الجليد
في آخر \num{700000} سنة يملك معظم قدرته هناك. (أ) بيّن أنه إذا كان التأثير $A_1\sin
\omega_1 t + A_2\sin\omega_2 t$، فإن نموذج التأخّر لا يستجيب إلا عند $\omega_1$ و$\omega_2$. (ب) بيّن
أن استجابة تحتوي على الحدّ $V^2$، مثل $\mathrm dV/\mathrm dt = (V_{\mathrm{eq}} -
V)/\tau + kV^2$، يمكن أن تُنشئ التردّدين $\omega_1 - \omega_2$ و$\omega_1 + \omega_2$، وذلك
بتعويض الحلّ الخطّي في $V^2$. (ج) حدّان من حدود المبادرة دوراهما \num{23.7}
و\num{19.0} ألف سنة يتداخلان فيعطيان نبضًا بأيّ دور؟ ماذا تستنتج بشأن الدورات ذات
\num{100000} سنة؟
\end{exercise}
% ledger: la04:precession, barker21:terminations

\begin{exercise}[$\star\star\star$]\label{exo:b2:orbital-forcing:11}
عند القطب الشمالي، الطاقة المتلقّاة لكلّ متر مربع خلال النهار القطبي هي $E = \int S_0(a/r)^2
\sin\delta\,\mathrm dt$ على مدى نصف السنة الذي تكون فيه $\delta > 0$. (أ) بيّن باستعمال برهان
\cref{prop:b2:orbital-forcing:annual} أن $E = S_0T\sin\varepsilon/(\pi\sqrt{1 - e^2})$. (ب) احسب
$E$ بوحدة \unit{GJ.m^{-2}} للقيمتين $\varepsilon = \ang{22.1}$ و\ang{24.5}، مع $e = 0.0167$ و$T$ =
365.25 يومًا. (ج) المبادرة لا تظهر في $E$. فلماذا يمكنها مع ذلك أن تغيّر
ذوبان غطاء جليدي؟
\end{exercise}
% ledger: nasa:milank-detail

\begin{exercise}[$\star\star\star$]\label{exo:b2:orbital-forcing:12}
(أ) احسب التأثير الإشعاعي لارتفاع ثنائي أكسيد الكربون من \qty{189.7}{ppm} إلى
\qty{279.4}{ppm}، مع \qty{3.93}{W.m^{-2}} لكلّ مضاعفة وقانون لوغاريتمي. (ب) المديان
المعطيان للذروة الجليدية الأخيرة ولآخر فترة بين جليدية هما من 188 إلى \qty{194}{ppm} ومن 266 إلى
\qty{282}{ppm}: احسب أصغر تأثير وأكبر تأثير يسمحان بهما. (ج) قارن بتغيّر
الإشماس العالمي السنوي في \cref{exo:b2:orbital-forcing:8} وبتأرجح
الإشماس الصيفي عند \ang{65} شمالًا. أيّهما الزناد وأيّهما المضخِّم؟
\end{exercise}
% ledger: data:b4-proxies-co2, ipcc:fig2-34, ipcc:erf-2xco2, ipcc:log-co2, data:b4-orbital

\begin{problem}[{مسألة نهاية الأسبوع --- ضوء الشمس عند \ang{65} شمالًا}]\label{pb:b2:orbital-forcing:1}
انتهت آخر فترة بين جليدية قبل نحو \num{116000} سنة. فكم تلقّى صيف دائرة
عرض الأغطية الجليدية الشمالية الكبرى من ضوء الشمس آنذاك، مقارنةً باليوم؟ استعمل معطيات
التمارين. تعطي الجداول، قبل \num{116000} سنة، $e = 0.0440$، و$\varpi = \ang{92.7}$،
و$\varepsilon = \ang{22.52}$؛ واليوم، بالتقريب نفسه، $e = 0.0172$، و$\varpi = \ang{101.4}$،
و$\varepsilon = \ang{23.45}$.

\medskip\noindent\textbf{الجزء الأول --- المدار آنذاك والآن.}
\begin{enumerate}
  \item احسب \omterm{def:b2:orbital-forcing:elements}{المبادرة المناخية} آنذاك والآن.
  \item احسب مسافتي الحضيض والأوج، بوحدات $a$، آنذاك والآن.
  \item عند أيّ طول للشمس كانت الأرض في الحضيض آنذاك، والآن؟ قرب أيّ انقلاب؟
  \item بأيّ عامل كان ضوء الشمس عند الحضيض أقوى منه عند الأوج، آنذاك والآن؟
  \item هل كان ميل المحور أكبر من اليوم أم أصغر؟ وماذا يعني ذلك لصيف
        العروض العليا؟
  \item احسب الإشماس العالمي السنوي المتوسط آنذاك والآن. هل الفرق معتبر؟
\end{enumerate}

\medskip\noindent\textbf{الجزء الثاني --- المسافة عند انقلاب يونيو.}
\begin{enumerate}[resume]
  \item بيّن انطلاقًا من \cref{prop:b2:orbital-forcing:distance} أنه عند انقلاب يونيو $r/a = (1 -
        e^2)/(1 - e\sin\varpi)$.
  \item احسب $r/a$ عند انقلاب يونيو آنذاك والآن.
  \item احسب $(a/r)^2$ آنذاك والآن.
  \item قارن بالتقريب $1 - 2e\sin\varpi$.
  \item بأيّ نسبة مئوية كان ضوء الشمس عند الانقلاب أضعف منه عند المسافة المتوسطة، آنذاك والآن؟
\end{enumerate}

\medskip\noindent\textbf{الجزء الثالث --- ضوء الشمس عند \ang{65} شمالًا.}
\begin{enumerate}[resume]
  \item أعطِ ميل الشمس عند انقلاب يونيو آنذاك والآن.
  \item احسب الزاوية الساعية للغروب عند \ang{65} شمالًا في ذلك اليوم، آنذاك والآن.
  \item استنتج طول النهار.
  \item احسب الإشماس اليومي المتوسط عند المسافة المتوسطة.
  \item احسب الإشماس اليومي المتوسط، آنذاك والآن.
  \item تعطي الجداول، لمنتصف يوليو، 397.60 و\qty{426.76}{W.m^{-2}}. لماذا تكون قيم
        الانقلاب أعلى؟
  \item أيّهما كان أثقل وزنًا في الفرق بين آنذاك والآن، ميل المحور أم المسافة؟
\end{enumerate}

\medskip\noindent\textbf{الجزء الرابع --- جواب الغطاء الجليدي.}
\begin{enumerate}[resume]
  \item في نموذج التأخّر مع $\tau = \qty{10}{kyr}$، ما الكسر الذي يكون الغطاء الجليدي قد أنجزه من تغيّر مفاجئ في
        $V_{\mathrm{eq}}$ بعد 5 و10 و\qty{20}{kyr}؟
  \item هبط إشماس يونيو عند \ang{65} شمالًا من قيمة عظمى قبل نحو \num{128000} سنة إلى
        القيمة الدنيا قبل \num{116000} سنة، أي نصف دورة مبادرة. ولدالة جيبية دورها
        \qty{23}{kyr}، ما الكسر الذي يتابعه الغطاء الجليدي من التأثير، وبأيّ تأخّر؟
  \item لماذا لا يستطيع غطاء جليدي أن يبلغ، خلال دورة مبادرة واحدة، الحجم الذي يكفي إشماس صيفي
        منخفض لإبقائه؟
  \item استغرق الجليد نحو \num{90000} سنة ليتراكم حتى الذروة الجليدية الأخيرة ونحو
        \num{10000} سنة ليذوب. لماذا لا يستطيع نموذج التأخّر أن ينتج هذا اللاتناظر؟
  \item كان ثنائي أكسيد الكربون بين 266 و\qty{282}{ppm} في آخر فترة بين جليدية وبين
        188 و\qty{194}{ppm} في الذروة الجليدية الأخيرة. كيف يساعد على أن يصبح تغيّر إقليمي في
        ضوء الشمس تغيّرًا عالميًا في المناخ؟
  \item خلاصة القول: ما الإشماس اليومي المتوسط عند \ang{65} شمالًا يوم انقلاب يونيو حين
        انتهت آخر فترة بين جليدية، قبل \num{116000} سنة؟
\end{enumerate}
\end{problem}
% ledger: berger91:rows, berger91:readme, data:b4-orbital, ipcc:ref-periods, ipcc:fig2-34, barker21:terminations
